\section{The logarithmic family and fixed-order closure}
\label{inf:sec-logarithmic}

\begin{definition}[The logarithmic input family]\label{inf:log-family}
With $m_0$ as in Definition~\ref{inf:input-families}, put
\[
 \mathcal P_{\log}=\left\{p\in\tfrac12\Z:p>1,\quad
 \exists m\ge m_0:\ |p-\widehat p_m|\le\frac{\log p}{8p}\right\}.
\]
The root $R$ and detuning are obtained by continuation from this
crossing. For sufficiently large members the crossing is unique,
and \eqref{inf:log-detuning-relation} gives $0<\tau\le\log p$.
Indeed $16\delta_0<8$, and non-coincidence excludes $p=\widehat p_m$.
Every fixed-width family $\mathcal P_{\mathrm{fix}}(C_{\mathrm{ap}})$ is contained
in $\mathcal P_{\log}$ from some point on.
\end{definition}

\begin{lemma}[Multipliers giving unbounded detuning]\label{inf:log-multiplier}
In each parity, $\mathcal P_{\log}$ has an unbounded selection with
$\widehat p_m|p-\widehat p_m|\to\infty$ and $\tau\to\infty$.
\end{lemma}
\begin{proof}
The congruence and shifted congruence lemmas give infinitely many
pairs $(P,Q)$ with $Q\equiv1\pmod4$, either desired parity of $P$,
and $0<e=Q|Q\alpha_{\mathrm{ar}}-P|<9/2$.
For each sufficiently large $Q$ choose the largest positive
$l\equiv1\pmod4$ with
\[
 l^2e\le\tfrac1{10}\log(lQ).
\]
There is such a multiplier eventually, and the set is finite since
$l^2/\log(lQ)\to\infty$ at fixed $Q$. Moreover $l\to\infty$:
every fixed admissible multiplier eventually satisfies the inequality
because $e<9/2$. Failure at $l+4$ gives the explicit two-sided bound
\begin{equation}\label{inf:multiplier-two-sided}
 \left(\frac l{l+4}\right)^2\frac{\log(lQ)}{10}
 <l^2e\le\frac{\log(lQ)}{10}.
\end{equation}
Set $P'=lP$, $Q'=lQ$, $p'=(P'-3)/2$ and let $p_0'$ be the
crossing with index $(Q'-1)/4$. Then $Q'\equiv1\pmod4$ and the
parity of $2p'$ is preserved. The crossing expansion gives
\[
 p_0'(p'-p_0')
 =-\frac{\alpha_{\mathrm{ar}}}{4}l^2Q(Q\alpha_{\mathrm{ar}}-P)
          +d_{\mathrm{ar}}+o(1).
\]
Since $l/(l+4)\ge1/5$, it follows that
\[
 \frac{\alpha_{\mathrm{ar}}}{1000}\log Q'-C
 \le p_0'|p'-p_0'|
 \le\frac{\alpha_{\mathrm{ar}}}{40}\log Q'+C.
\]
Here $p'/Q'\to\alpha_{\mathrm{ar}}/2$ and
$p_0'/p'\to1$. As $\alpha_{\mathrm{ar}}/40<1/8$, the upper bound
puts these inputs in $\mathcal P_{\log}$ eventually; the lower
bound tends to infinity. Equation~\eqref{inf:log-detuning-relation}
then gives $\tau\to\infty$. Successive primary crossings are
separated by $2\alpha_{\mathrm{ar}}+o(1)$, so another crossing
cannot place these same inputs in a fixed-width family.
\end{proof}

\begin{theorem}[Newton closure at order twenty]\label{inf:fixed-order-newton}
Fix $5/3<q_{\mathrm{gap}}<2$ and
$0<\xi_{\mathrm{inv}}<2-q_{\mathrm{gap}}$, a compact split interval,
and a positive gap constant. At every sufficiently large logarithmic
input whose quartic domain has gap at least
$\gamma=\kappa_{\mathrm{gap}}\tau p^{-q_{\mathrm{gap}}}$,
the order-$20$ centre has an exact real zero within
\[
 \nrm{z_*-z_0}_{\Xc_p}\le\tau p^{-11}.
\]
The zero gives a bounded strictly convex non-ball analytic domain
with the frequency-one Schiffer data.
\end{theorem}
\begin{proof}
Fix the parameters in the statement. Choose
$\alpha_{\mathrm{loc}}>0$ and an internal loss smaller than
$\xi_{\mathrm{inv}}$ as in Theorem~\ref{inf:mixed-conditioning}.
Fix the angular window constant and then $K_{\mathrm{core}}$ to
control the global Green sums. Fix the radial head cutoff.
Take $N=20$ and fix the complex angular disc and coefficient moments
needed by the centre, clamping and material estimates.
These choices all precede the threshold in $p$.

Lemma~\ref{inf:detuning-dependencies} supplies the true-centre
hypothesis of Theorem~\ref{inf:mixed-conditioning} and the
same-space residual and material factorization. Its row 12 supplies
the sharp material and nonlinear bounds uniformly in this range.
They give
\[
 \nrm{\AN}\le\tau^{-1}p^{B_{\log}+o(1)},\qquad
 B_{\log}=\tfrac92+q_{\mathrm{gap}}+\xi_{\mathrm{inv}}<\tfrac{13}2,
 \qquad \sup\nrm{D^2\Fnl}\le p^{3+o(1)},
\]
where $\AN=\Qmat^{-1}J_D^{-1}$ is a bounded isomorphism.
On the closed ball of radius $\mathfrak r_p=\tau p^{-11}$,
the three contraction quantities are bounded by
\begin{align}
 \nrm{\AN\Fnl(z_0)}/\mathfrak r_p
 &\le p^{B_{\log}-7+o(1)}\le p^{-1/2+o(1)},
 \label{inf:log-gate-one}\\
 \nrm{I-\AN D\Fnl(z_0)}
 &\le\tau p^{B_{\log}-17+o(1)}
       \le(\log p)p^{-21/2+o(1)},\label{inf:log-gate-two}\\
 \sup\nrm{\AN D^2\Fnl}\,\mathfrak r_p
 &\le p^{B_{\log}-8+o(1)}\le p^{-3/2+o(1)}.
 \label{inf:log-gate-three}
\end{align}
All tend to zero uniformly in $0<\tau\le\log p$.
The source displacement is below one and the total shape norm
below $1/4$, so the entire closed ball is in the working region.
For large $p$ the map $z\mapsto z-\AN\Fnl(z)$ maps that ball
to itself and has Lipschitz constant at most $1/2$, by the same
Taylor integral inequalities used in Theorem~\ref{inf:newton-threshold}.
Completeness and injectivity of $\AN$ give an exact real zero.

Its normalized $C^2$ shape is $O(\tau p^{-2})+O(\tau p^{-11})$,
which tends to zero. The physical seed coefficient is
$-\theta/(16r(1-r)p)+O(\tau p^{-2}\mathfrak P_{20}(\log p))$;
the correction changes it by at most $C\tau p^{-10}$.
Division by $\tau/p$ makes both errors tend to zero.
Thus the curvature and non-ball arguments in
Section~\ref{inf:sec-geometry} apply. The completed-space
identification in Proposition~\ref{inf:classical-solution} supplies
the classical field and its two boundary traces in the integer
dimension $2p$.
\end{proof}

\begin{proposition}[Assembly of the logarithmic family]\label{inf:log-assembly}
For every fixed $0<\varepsilon_{\mathrm{sp}}<1/4$ and $0<\eta_{\mathrm{cnt}}<1/3$,
every sufficiently large $p\in\mathcal P_{\log}$ has the domains of
Theorem~\ref{inf:fixed-order-newton} at all but
$O_{\mathrm{sp},\eta_{\mathrm{cnt}}}(p^{2/3+\eta_{\mathrm{cnt}}})$ interior integer splits.
Distinct retained splits give pairwise non-similar domains.
The similarity-class lower limit is at least $1/2$ after division
by $n=2p$, along this family's dimensions.
\end{proposition}
\begin{proof}
Choose
\[
 q_{\mathrm{gap}}=2-\eta_{\mathrm{cnt}}/2,\qquad
 \xi_{\mathrm{inv}}=\eta_{\mathrm{cnt}}/4,\qquad
 \alpha_{\mathrm{loc}}=\eta_{\mathrm{cnt}}/16.
\]
These obey all strict inequalities in the conditioning argument.
Definition~\ref{inf:log-family} and
Lemma~\ref{inf:detuning-dependencies} supply its arithmetic,
low-state and uniform-constant hypotheses.
Corollary~\ref{inf:log-flow-count} gives the selected quartic gaps
and at most $C\sqrt{\log p}\,p^{2/3+\eta_{\mathrm{cnt}}/2}
=O(p^{2/3+\eta_{\mathrm{cnt}}})$ exceptions.
For every retained split apply Theorem~\ref{inf:fixed-order-newton}
at the single fixed accuracy twenty. All its thresholds are
uniform over the compact split interval. Lemma~\ref{inf:similarity-classes}
proves non-similarity and the lower-limit conclusion by the same
integer count and then by letting the fixed margin tend to zero.
\end{proof}
